\answer{ex:11:1}{(a)~$T=k^2/\bigl(\Phi(\Delta I/I)^2\bigr)=9\,\text{s}$. (b)~$45$ sensores, mancha de $\approx7\times7$, resolução $\approx7$ vezes pior.}
\answer{ex:11:2}{(a)~$2$--$3.3$ bits por disparo, $22$--$34\,\%$ do limite ($9.1$--$9.8$ bits). (b)~$\log_2\binom{400}{20}\approx111$ bits; em média $1$ tipo em comum.}
\answer{ex:11:3}{(a)~$\sigma/9\le I\le9\sigma$, $1.9$ ordem de grandeza. (b)~$6$ e $7$.}
\answer{ex:11:4}{(a)~$f<343/0.44\approx780$~Hz. (b)~$625$ disparos, $125$ neurônios.}
\answer{ex:11:5}{$21$ bits por evento, $4.8\cdot10^5$ eventos/s, $7$ eventos por pixel de borda: $136$ pixels/s. Uma câmera comum dá $1.25$ quadro/s.}
\answer{ex:11:6}{Em logaritmos, $\approx1.0\,\text{s}$ nos dois sentidos (teoria $0.96\tau$); linearmente, $0.2\,\text{s}$ para cima e $3.0\,\text{s}$ para baixo (teoria $0.18\tau$ e $3.0\tau$).}
\answer{ex:11:7}{(a)~$\ln I$ tem distribuição logística com mediana $\ln\sigma$ e escala $1$: dispersão $\pi/\sqrt3$ em logaritmos naturais, $0.79$ ordem de grandeza. (b)~$r=r_{\max}\Phi_I(I)^2$.}
